\documentclass[10pt,a4paper]{amsart}
\usepackage[margin=19mm]{geometry}
\usepackage{amsmath,amssymb,mathtools}
\usepackage[scr=rsfs,cal=euler]{mathalfa}
\usepackage{xfrac}
\usepackage[unicode,hidelinks,bookmarks=false]{hyperref}
\newcommand{\ie}{i.e.\ }
\newcommand{\cf}{cf.\ }
\newcommand{\resp}{respectively\ }
\newcommand{\Subsection}{\subsection}
\providecommand{\nobreakdash}{\nobreak\hbox{-}}
\setlength{\emergencystretch}{2em}
\multlinegap0pt
\usepackage{amssymb}
\usepackage{mathtools}
\usepackage{tikz}
\newtheorem{theorem}{Theorem}
\newtheorem{corollary}{Corollary}
\newcommand{\bry}[1]{\partial #1}
\newcommand{\cubical}[2]{Q_{#1}^{#2}}
\newcommand{\slice}[2]{S_{#1,#2} } 
\title{Overlap-Helly theorems}
\author{Andreas F. Holmsen, Alfredo Hubard}
\date{Source: arXiv:2609.10023v1 (CC BY 4.0)}
\begin{document}
\maketitle

\noindent\emph{Source and licence.} Overlap-Helly theorems. Version arXiv:2609.10023v1 (CC BY 4.0), distributed by arXiv under the Creative Commons Attribution 4.0 International licence, http://creativecommons.org/licenses/by/4.0/, as recorded in the arXiv licence metadata for this exact version and shown on the abstract page https://arxiv.org/abs/2609.10023v1. Full licence notice: \texttt{licenses/arxiv-2609.10023-CC-BY-4.0.txt}.\par\smallskip

\noindent\textit{Editorial scope.} Counted unit: the article's theorem version (source0.tex lines 152--154). The article's complete original proof of it is delivered with it (source0.tex lines 337--355, one \texttt{proof} environment of 19 lines, which the article places away from the statement). No mathematics is written, repaired or re-derived: the statement and the proof are the article's own lines, in the article's own notation, and no retained line is substituted or re-flowed.  Retained uncounted as the source-local context the proof needs: lines 174--175; lines 201--203.  Nothing was omitted from any retained span, so the package carries no omission record.  No retained line contains a citation, so the wrapper carries no bibliography and no bibliography entry.  No retained line contains a figure environment, an image-inclusion command or drawing code, so no image is delivered and the package declares no asset dependency.  Editorial formatting only, in addition: an AMS \texttt{amsart} preamble that restates the article's own declarations and macros used by the retained lines, namely the environments \texttt{theorem}, \texttt{corollary} on separate counters, together with the standard packages those lines need.  The retained environments print in this document as Theorem 1, Corollary 1 in order of appearance, whereas the article prints the same items with the numbering of its own full document.  The complete native source of the article as distributed in the arXiv e-print for this exact version is bundled unchanged as \texttt{sources/209868/source0.tex}.
\begin{corollary}\label{cor:good covers}
Let $Y$ be a $d$-dimensional simplicial complex such that $\pi_d(Y)$ vanishes. Let $F_1$, $\dots$, $F_{d+1}$ be finite families of subcomplexes of $Y$ such that for every nonempty $I\subset [d+1]$ and for every choice $S_i\in F_i$ with $i\in I$, we have $\pi_{d-|I|}(\bigcap_{i\in I}S_i)$ is trivial. Then there is a point $y\in \|Y\|$ that belongs to every member of one of the families $F_i$.  \end{corollary}

\begin{theorem}\label{t:top col helly}
Let $Y$ be a $d$-dimensional simplicial complex. For every continuous map  $f \colon \cubical{n}{d+1} \to \|Y\|$ there exists an $i \in [d+1]$ and a point $y\in \|Y\|$ such that $f^{-1}(y)$ intersects $\slice{i}{j}$ for every $1\leq j\leq n$.
\end{theorem}

\begin{theorem}\label{thm:convex version}
Let $m >  d\geq 1$ be integers, let $F_1, \dots, F_{d+1}$ be finite families of compact convex sets in $\mathbb{R}^m$, and let $Y$ be a $d$-dimensional simplicial complex. If $C_1\cap \cdots \cap C_{d+1}\neq\emptyset$ for every choice $C_1\in F_1, \dots, C_{d+1}\in F_{d+1}$, then for every continuous map $f \colon \mathbb{R}^m \to \|Y\|$ there exists a point $y\in \|Y\|$ such that $f^{-1}(y)$ intersects every member of one of the families $F_i$.
\end{theorem}

\begin{proof}[Proof of Corollary \ref{cor:good covers}]
Without loss of generality we may assume that $|F_i| = n$ for every $i \in [d+1]$ by adding copies of some of the sets if necessary. Now arbitrarily label the members of $F_i$ as 
$C_{i,1}, \dots, C_{i,n}$. We will construct a continuous map $f \colon \cubical{n}{d+1} \to \|Y\|$ such that 
$f(\slice{i}{j}) \subset C_{i,j}$
for every $i\in [d+1]$ and $j\in [n]$. The result then follows from Theorem \ref{t:top col helly} (which is a special case of Theorem \ref{thm:convex version}).

It remains to construct the map $f$. We consider the natural cubical cell structure of $\cubical{n}{d+1}$, define the map on its vertices, and proceed inductively. For a vertex $v = (t_1, \dots, t_{d+1}) \in [n]^{d+1}$ we choose (arbitrarily) a point 
\[p \in \bigcap_{i\in [d+1]} C_{i,t_i}, \]
which is possible since this intersection is nonempty: it is the case $I = [d+1]$ of the hypothesis, where $\pi_{d-|I|} = \pi_{-1}$ is trivial precisely when $\bigcap_{i\in [d+1]} C_{i,t_i}$ is nonempty. We set $f(v) = p$. This defines the map $f$ on the 0-skeleton of $\cubical{n}{d+1}$. 
Now suppose $f$ has been defined on the $k$-skeleton for some $0\leq k < d$ such that 
\[\tau \in \slice{i}{j} \implies f(\tau) \subset C_{i,j}\] for every cell $\tau$ in the $k$-skeleton of $\cubical{n}{d+1}$.
Now consider a $(k+1)$-cell $\sigma \in \cubical{n}{d+1}$. The cell $\sigma$ is contained in the intersection of a unique $(d-k)$-tuple of axis parallel hyperplanes, that is, there is a $(d-k)$-element subset $I\subset [d+1]$ and a set of integers $\{t_i\}_{i\in I}$, with $t_i \in [n]$, such that 
\[\sigma \subset \bigcap_{i\in I} \slice{i}{t_i}.\]
Since the boundary $\bry{\sigma}$ lies in the $k$-skeleton,  we have 
\[f\big( \bry{\sigma} \big) \subset  \bigcap_{i\in I} C_{i,t_i}. \]
By assumption,
$\pi_{k}(\bigcap_{i\in I}C_{i,t_i})$ is trivial, 
and since $\bry{\sigma}$ is homeomorphic to a $k$-sphere, we can extend $f$ on $\sigma$ in such a way  that $f(\sigma) \subset \bigcap_{i\in I} C_{i,t_i}$. Proceeding in this way defines the map on the $d$-skeleton of $\cubical{n}{d+1}$ such that $f(\slice{i}{j})\subset C_{i,j}$ for every $i\in [d+1]$ and $j\in [n]$. Finally, we extend $f$ to the rest of $\cubical{n}{d+1}$ by extending over each individual $(d+1)$-cell. Such a cell is contained in no slice, so this is the case $I = \emptyset$, where $\bigcap_{i\in \emptyset} C_{i,t_i} = \|Y\|$ and there is no further containment requirement to be satisfied; the extension exists because $\pi_d(Y)$ is trivial, which is the $I=\emptyset$ instance of the hypothesis.
\end{proof}

\end{document}
